\section{Estimating $\pi$}

\begin{center}
    % \renewcommand{\arraystretch}{1.5}
    \begin{tabular}{|l|p{10cm}|}
        \hline
        \textbf{Date} & \textbf{Description}                                 \\
        \hline
        May 23, 2026  & Created cargo project, implemented pi estimation     \\
        \hline
        May 24, 2026  & Created unit tests for pi estimation, created \LaTeX
        logbook                                                              \\
        \hline
    \end{tabular}
\end{center}

\subsection{Goal}

Calculate an estimate of $\pi$ to at least an accuracy of $3.14$.

\subsection{Simulation}

We will generate $\pi$ from its definition in the area of a circle:
$A = \pi r^2$. Suppose we inscribe the circle in a square of side length $2r$.
Let $p$ be the ratio of the area of the circle to the area of the square. Then
we have

\[
    p = \frac{\pi r^2}{(2r)^2} = \frac{\pi}{4} \implies \pi = 4p.
\]

Indeed, we can estimate $p$ by randomly (uniformly) generating points in the
square, and counting the fraction of points that fall inside the circle.

Suppose we center the circle and the square at the origin, and fix the square to
be aligned with the axis. Then we can generate points $(X_i, Y_i)$ where
$X_i, Y_i \sim \text{Uniform}(-r, r)$. We can check if a point is inside the
circle by checking if $X_i^2 + Y_i^2 \leq r^2$.

Let $I_i$ be the indicator variable for the event that the $i$-th point is
inside the circle. Then we can estimate $p$ as

\[
    \hat{p} = \frac{1}{n} \sum_{i=1}^n I_i.
\]

where $n$ is the total number of points generated. Finally, we estimate $\pi$ as

\[
    \hat{\pi} = 4 \hat{p} = \frac{4}{n} \sum_{i=1}^n I_i.
\]


\begin{note}{Reducing Scope}{}
    Note that there is a rotational symmetry in the problem, so we can just
    narrow the point generation to the first quadrant
    ($X_i, Y_i \sim \text{Uniform}(0, r)$). The rest of the process is the same.
\end{note}

\subsection{Rust Implementation}

The Rust implementation of the above is not particularly complex, just a simple
loop to generate and count points.
